% Generated draft: requires mathematical and source review.
\documentclass[11pt]{article}
\usepackage[margin=25mm]{geometry}
\usepackage{fontspec}
\setmainfont{Latin Modern Roman}
\usepackage{amsmath,amssymb,mathtools,unicode-math}
\setmathfont{Latin Modern Math}
\usepackage{enumitem,needspace,xurl,hyperref}
\setlength{\parindent}{0pt}
\setlength{\parskip}{5pt}
\setlength{\emergencystretch}{3em}
\newcommand{\N}{\mathbb{N}}
\newcommand{\Z}{\mathbb{Z}}
\newcommand{\Q}{\mathbb{Q}}
\newcommand{\R}{\mathbb{R}}
\newcommand{\C}{\mathbb{C}}
\newcommand{\F}{\mathbb{F}}
\newcommand{\A}{\mathbb{A}}
\newcommand{\PP}{\mathbb{P}}
\newcommand{\E}{\mathbb{E}}
\newcommand{\eps}{\varepsilon}
\newcommand{\dd}{\,\mathrm{d}}
\begin{document}
\section*{UnsolvedMath ID: 30006161}
\textbf{Source problem code:} catalogue entry
\subsection{Definitions and mathematical statement}
Let \(S^2=\{x\in\R^3:\|x\|=1\}\) with its usual topology, and let \(\mathbb{RP}^2=S^2/(x\sim -x)\) with the quotient topology. For each \(X\in\{S^2,\mathbb{RP}^2\}\), let \(\mathcal K(X)\) be the hyperspace of nonempty closed subsets of \(X\). Equip it with the Vietoris topology: for open subsets \(U_1,\ldots,U_n\) of \(X\), where \(n\geq 1\), the sets
\[
\langle U_1,\ldots,U_n\rangle
=\left\{K\in\mathcal K(X):K\subseteq\bigcup_{i=1}^nU_i\ \text{and}\ K\cap U_i\ne\varnothing\ \text{for every }i\right\}
\]
form a basis. Write
\[
\mathcal C(X)=\{K\in\mathcal K(X):K\text{ is connected}\}.
\]
A chain is a family \(\mathscr C\subseteq\mathcal C(X)\) such that, for all \(A,B\in\mathscr C\), either \(A\subseteq B\) or \(B\subseteq A\). It is maximal if no strictly larger subfamily of \(\mathcal C(X)\) is a chain. Let \(\Phi(X)\) denote the set of these maximal chains. Each maximal chain is a nonempty closed subset of \(\mathcal K(X)\); equip \(\Phi(X)\) with the subspace topology inherited from the Vietoris hyperspace \(\mathcal K(\mathcal K(X))\).

Let \(\operatorname{Homeo}(X)\) be the group of all homeomorphisms from \(X\) onto itself. Its action on \(\Phi(X)\) is
\[
h\cdot\mathscr C=\{h[K]:K\in\mathscr C\}.
\]
A subset of \(\Phi(X)\) is comeager if its complement is a countable union of nowhere dense subsets of \(\Phi(X)\). A generic chain means a chain whose orbit under this action is comeager.

The two questions, asked separately, are whether
\[
\exists\,\mathscr C\in\Phi(S^2)\quad
\operatorname{Homeo}(S^2)\cdot\mathscr C
\text{ is comeager in }\Phi(S^2),
\]
and whether
\[
\exists\,\mathscr C\in\Phi(\mathbb{RP}^2)\quad
\operatorname{Homeo}(\mathbb{RP}^2)\cdot\mathscr C
\text{ is comeager in }\Phi(\mathbb{RP}^2).
\]
The hyperspace construction above spells out the standard interpretation of the supplied phrase ``natural compact hyperspace topology.'' The supplied excerpt does not explicitly define that topology or specify whether the empty set is included among connected closed subsets. Here the nonempty convention is explicit; if the empty set is allowed, every maximal chain contains it, and adjoining or deleting that common least member gives the corresponding formulation. These conventions require confirmation against the full source context. The supplied record labels both questions open; that label is supplied curation, not an independently verified conclusion.
\subsection{Short English statement}
For the sphere and, separately, the real projective plane, ask whether there is a maximal family of nested connected closed sets whose images under all homeomorphisms form a comeager subset of the space of maximal chains. The draft makes explicit the usual Vietoris topology and nonempty-set convention, which the supplied excerpt leaves implicit.
\subsection{Sources}
\begin{itemize}
\item Local UnsolvedMath registry: supplied complete source record, permanent ID \(30006161\), title ``Generic Maximal Chains on Exceptional Surfaces,'' exact source problem code \texttt{catalogue entry}. The record gives proposed year \(2025\) and supplies no individual proposer.
\item Supplied primary report citation: \emph{Set Theory} (2025), Oberwolfach Reports, pp.\ 75--140, DOI \texttt{10.4171/owr/2025/2}. \url{https://doi.org/10.4171/owr/2025/2}. The supplied passage identifies this as Question 3 and refers to ``[1], Question 1.3''; the bibliographic identity of [1] is not supplied.
\item Supplied literature source: ``Oberwolfach Report 2/2025.'' \url{https://content.ems.press/assets/public/full-texts/serials/owr/22/1/14299082/online/10.4171-owr-2025-2.pdf}.
\item Supplied literature assessment, checked \texttt{2026-08-22}: the primary report labels the generic-chain questions for the sphere and real projective plane as still open, and the supplied curation reports that no later primary resolution was located by that date. This is supplied dated curation, not independently verified here; time-sensitive status and cited claims require independent verification.
\item The registry assigns statement status \texttt{corrected\_verified} and states that the formulation was reconstructed and checked against the official Oberwolfach report. This is a supplied source-verification claim, not verification performed for this draft.
\end{itemize}
\end{document}
